\chapter{QS-Maßnahmen}
Im Folgenden werden die durchgeführten Maßnahmen zur Qualitätssicherung (QS) der Studienarbeit beschrieben.

Farberklärung:
\begin{itemize}
    \item {\textcolor{yellow!80!brown}{QS-Maßnahme in Arbeit}}
    \item {\textcolor{ForestGreen}{QS-Maßnahme erfolgreich und abgeschlossen}}
    \item {\textcolor{red!80!black}{QS-Maßnahme nicht erfolgreich und abgeschlossen}}
\end{itemize}
\section{Besprechung mit Betreuer}
Regelmäßige Besprechungen mit dem Betreuer der Studienarbeit wurden durchgeführt, um den Fortschritt zu überprüfen, Feedback zu erhalten und sicherzustellen, dass die Arbeit den Anforderungen entspricht. Diese Besprechungen fanden in Form von Online-Meetings und persönlichen Treffen statt.
Bisherige Treffen: 
\begin{itemize}
    \item 28.07.2025: Erstes Treffen zur Besprechung des Themas und der Zielsetzung der Studienarbeit.
    \item 09.10.2025: Treffen an der DHBW zur Besprechung des Projektplans und der ersten Schritte. Übergabe der bestellten Hardware.
    \item 13.112025: Besprechung des defekten Pis und der weiteren Vorgehensweise.
\end{itemize}
Teilnehmende: Christoph Niederer, Luca Riebel, Silvan Müller.

\textcolor{yellow!80!brown}{\textbf{QS-Maßnahme in Arbeit}}
\section{Testen der Hardware}
Während der Erstellung des ersten Konzeptes (05.11.2025) wurde der Raspberry Pi getestet, indem der erste Software-Prototyp darauf ausgeführt wurde. Dabei wurde festgestellt, dass der Raspberry Pi einen defekt aufweist, da er die SD Karte nicht auslesen kann. Dies wurde dem Betreuer mitgeteilt, um eine Lösung zu finden.

Durch diese frühzeitige QS-Maßnahme, das testen der Hardware, konnte ein potenzielles Problem erkannt und adressiert werden, bevor es den Fortschritt der Studienarbeit erheblich beeinträchtigt hätte.


17.11.2025: Der Raspberry Pi wurde ersetzt und der Prototyp erfolgreich darauf getestet. Im nächsten Schritt wird die Integration mit dem Display und Spiegel getestet.
Teilnehmende: Christoph Niederer, Luca Riebel.

{\textcolor{ForestGreen}{QS-Maßnahme erfolgreich und abgeschlossen}}

\section{Feedback zur Gliederung}
Am 17.11.2025 wurde die Gliederung der Studienarbeit erstellt und soll nun mit dem Betreuer besprochen werden, um sicherzustellen, dass sie den Anforderungen entspricht und alle relevanten Themen abdeckt.

Teilnehmende: Christoph Niederer, Luca Riebel

{\textcolor{yellow!80!brown}{QS-Maßnahme in Arbeit}}

\section{Testen des Spiegels}
Am 19.11.2025 wurde der Spiegel in Kombination mit dem Raspberry Pi + Display getestet. Das frühzeitige Testen ist notwendig, damit erkannt werden kann, ob das Display hinter dem Spionspiegel gut sichtbar ist.

Teilnehmende: Christoph Niederer, Luca Riebel

{\textcolor{ForestGreen}{QS-Maßnahme erfolgreich und abgeschlossen}}
